% ============================================================
\chapter*{Preface}
\addcontentsline{toc}{chapter}{Preface}
% ============================================================

This course on \emph{Advanced Point-Set Topology} is aimed at Master's and doctoral
students in pure mathematics, as well as researchers in theoretical computer science
wishing to deepen their understanding of the topological foundations of domain theory.

\section*{Motivations}

General topology, as taught in undergraduate and early graduate programs, focuses almost
exclusively on Hausdorff (separated) spaces. While natural for analysis and differential
geometry, this restriction conceals a substantial portion of the theory. Indeed, many
fundamental structures in mathematics and computer science fail to satisfy the $T_2$
separation axiom:

\begin{itemize}
  \item The spectrum of a commutative ring, endowed with the Zariski topology, is
    generally a $T_0$-space that is not $T_1$.
  \item Scott domains, essential in denotational semantics, carry a $T_0$ topology
    that is not $T_1$.
  \item Locales, which generalize topological spaces by eliminating reference to points,
    provide a natural framework for constructive logic and topos theory.
  \item Quasi-metrics, where symmetry is abandoned, model asymmetric cost notions
    in computer science and optimization.
\end{itemize}

The goal of this course is to systematically develop topological theory in this broader
setting, emphasizing the interactions between topology, order theory, and theoretical
computer science.

\section*{Organization of the Course}

The course is divided into ten chapters, organized into three thematic parts:

\medskip
\noindent\textbf{Part~I: Foundations (Chapters~1--3).}
We begin with a systematic review of general topology enriched with categorical language
(Chapter~1), then study $T_0$-spaces via the specialization order, establishing the
fundamental bridge between topology and order theory (Chapter~2), before addressing sober
spaces and the theory of locales providing the ``pointless'' framework of topology
(Chapter~3).

\medskip
\noindent\textbf{Part~II: Asymmetric Structures (Chapters~4--6).}
We develop the theory of quasi-metrics and asymmetric topology (Chapter~4), bitopological
spaces capturing the duality between a topology and its conjugate (Chapter~5), and then
Scott topology on ordered sets and domains, fundamental for denotational semantics
(Chapter~6).

\medskip
\noindent\textbf{Part~III: Convergence, Compactness, and Applications (Chapters~7--10).}
We study generalized convergence spaces and filter theory (Chapter~7), compactness beyond
the Hausdorff setting with its multiple variants (Chapter~8), applications of asymmetric
topology to theoretical computer science (Chapter~9), and finally topological fixed point
theory with applications to domains and semantics (Chapter~10).

\section*{Prerequisites}

The reader is assumed to be familiar with:
\begin{itemize}
  \item Basic general topology: topological spaces, continuity, compactness, connectedness,
    separation axioms $T_0$--$T_4$ (advanced undergraduate/early graduate level).
  \item Elementary notions of ordered set theory: partial orders, lattices, completeness.
  \item The basics of categorical language: categories, functors, natural transformations,
    adjunctions.
  \item Elementary theory of metric spaces and convergence.
\end{itemize}

\section*{Conventions and Notation}

Throughout this course, we adopt the following conventions:
\begin{itemize}
  \item $(X,\tau)$ denotes a topological space, where $\tau$ is the topology (the collection
    of open sets). When the topology is clear from context, we write simply~$X$.
  \item $\mathcal{O}(X)$ denotes the lattice of open sets of~$X$.
  \item $\mathcal{C}(X)$ denotes the collection of closed sets of~$X$.
  \item $\overline{A}$, $\mathrm{Int}(A)$ denote the closure and interior of
    $A \subseteq X$, respectively.
  \item $\leq_X$ or $\sqsubseteq$ denotes the specialization order on a $T_0$-space.
  \item $\downarrow\! x = \{y : y \leq x\}$ and $\uparrow\! x = \{y : x \leq y\}$.
  \item For a subset $A$, we write $\downarrow\! A = \bigcup_{a\in A}\downarrow\! a$
    and $\uparrow\! A = \bigcup_{a\in A}\uparrow\! a$.
  \item $\mathbf{Top}$ is the category of topological spaces and continuous maps.
  \item $\mathbf{Top}_0$ is the full subcategory of $T_0$-spaces.
  \item $\mathbf{Sob}$ is the full subcategory of sober spaces.
  \item $\mathbf{Loc}$ is the category of locales.
  \item $\mathbf{Frm}$ is the category of frames.
  \item $\mathbf{Dcpo}$ is the category of dcpos and Scott-continuous functions.
  \item $\mathbf{Dom}$ denotes the category of domains (context-dependent).
  \item $\omega$ denotes the first infinite ordinal, identified with $\N$.
  \item $\mathbb{S}$ denotes the Sierpi\'nski space $\{0,1\}$ with topology
    $\{\emptyset, \{1\}, \{0,1\}\}$.
\end{itemize}

\section*{Pedagogical Approach}

Each chapter contains:
\begin{itemize}
  \item Precise \textbf{definitions} and \textbf{theorems} with complete proofs.
  \item Detailed \textbf{examples} and \textbf{counterexamples} illustrating the scope and
    limitations of results.
  \item \textbf{Exercises} of varying difficulty, some accompanied by hints.
  \item \textbf{Intuition} and \textbf{Warning} boxes to guide understanding.
\end{itemize}

\section*{References}

This course draws primarily on the following works:
\begin{itemize}
  \item \textsc{Gierz, Hofmann, Keimel, Lawson, Mislove, Scott},
    \emph{Continuous Lattices and Domains}, Cambridge University Press, 2003.
  \item \textsc{Goubault-Larrecq}, \emph{Non-Hausdorff Topology and Domain Theory},
    Cambridge University Press, 2013.
  \item \textsc{Johnstone}, \emph{Stone Spaces}, Cambridge University Press, 1982.
  \item \textsc{Vickers}, \emph{Topology via Logic}, Cambridge University Press, 1989.
  \item \textsc{Abramsky, Jung}, \emph{Domain Theory}, Handbook of Logic in Computer
    Science, vol.~3, Oxford University Press, 1994.
  \item \textsc{Kelly}, \emph{Bitopological Spaces}, Proc.\ London Math.\ Soc., 1963.
  \item \textsc{Fletcher, Lindgren}, \emph{Quasi-Uniform Spaces}, Marcel Dekker, 1982.
  \item \textsc{Smyth}, \emph{Topology}, Handbook of Logic in Computer Science, vol.~1,
    Oxford University Press, 1992.
  \item \textsc{Escard\'o}, \emph{Synthetic Topology of Data Types and Classical Spaces},
    ENTCS, 2004.
  \item \textsc{Nachbin}, \emph{Topology and Order}, Van Nostrand, 1965.
\end{itemize}

\section*{Acknowledgements}

This course has benefited from the teachings and discussions with many colleagues working
at the interface of topology, order theory, and theoretical computer science. We thank
in particular those students whose questions and remarks helped improve the presentation.

\bigskip
\hfill\emph{The author}
